\documentclass[11pt]{article}
\usepackage[T1]{fontenc}
\usepackage{textcomp}
\usepackage[margin=0.75in]{geometry}
\usepackage{amsmath,amssymb,amsthm}
\usepackage{array,booktabs}
\usepackage{hyperref}
\setlength{\parindent}{0pt}
\newtheorem{theorem}{Theorem}
\theoremstyle{definition}
\newtheorem{definition}{Definition}
\title{Flow Proof Artifact: prop-31-angle-in-segment-and-right-angle}
\author{Generated by \texttt{flow doc proof}}
\date{Flow Proof Book}
\begin{document}
\maketitle
In a circle the angle in a greater segment is less than a right angle, and in a less segment greater.\\[0.5em]
\textbf{Source.} euclid — Elements, Book III, Proposition 31\\[0.5em]
\bigskip
\noindent{\large \textbf{Derived fact 1.} \textit{In a circle the angle in a greater segment is less than a right angle, and in a less segment greater} \hfill the Euclidean plane \cdot Euclid Book III \cdot Proposition 31: the angle in a greater segment is less than a right angle}\par
\smallskip\noindent\textit{Coordinate: the Euclidean plane \cdot Euclid Book III \cdot Proposition 31: the angle in a greater segment is less than a right angle}\par
\smallskip\noindent\textit{Source: euclid — Elements, Book III, Proposition 31}\par
\smallskip\noindent\textit{Built on: proposition 20: the central angle is double the inscribed angle on the same arc, for Euclid Book III on the Euclidean plane, proposition 11: erect a perpendicular at a point on a line, for Book I of the Elements in on the Euclidean plane}\par
\medskip
\noindent\textbf{Goal.} In a circle the angle in a greater segment is less than a right angle, and in a less segment greater\par
\noindent$\displaystyle \text{segment angle compares to right}$\par
\medskip
\noindent
\begin{tabular}{@{} >{\raggedright\arraybackslash}p{0.06\textwidth} >{\raggedright\arraybackslash}p{0.47\textwidth} >{\raggedleft\arraybackslash}p{0.41\textwidth} @{}}
\textbf{\#} & \textbf{Proof} & \textbf{Mathematics} \\
\midrule
\textbf{1.} & We prove that proposition 31: the angle in a greater segment is less than a right angle for Euclid Book III the Euclidean plane. & \\[0.45em]
\textbf{2.} & Let seg = greater\_segment. & \\[0.45em]
\textbf{3.} & Let angle = inscribed\_angle. & \\[0.45em]
\textbf{4.} & From step 2 and step 3, we can deduce that angle < one right angle. & $\displaystyle angle < 90^\circ$ \\[0.45em]
\textbf{5.} & We invoke the derived fact governing Euclid Book III the Euclidean plane: proposition 20: the central angle is double the inscribed angle on the same arc, for Euclid Book III on the Euclidean plane. & \\[0.45em]
\textbf{6.} & We invoke the derived fact governing Book I of the Elements in the Euclidean plane: proposition 11: erect a perpendicular at a point on a line, for Book I of the Elements in on the Euclidean plane. & \\[0.45em]
\textbf{7.} & From step 5 and step 6, this implies segment angle compares to right. Hence proven. & $\displaystyle \text{segment angle compares to right}$ \\[0.45em]
\bottomrule
\end{tabular}
\label{thm:Geometry-Euclid Book III-proposition_31_the_angle_in_a_greater_segment_is_less_than_a_right_angle}
\par\medskip\hrule
\end{document}
